%Équilibres binaires (S-L) 13, Faynot


\setcounter{equation}{0}
\setcounter{Q}{0}

Toutes les expériences et mesures, décrites dans cet exercice, sont menées sous une pression constante $P = P_{atm} = 1.01 ~ \si{\bar}$. Les courbes de refroidissement de mélanges cadmium-plomb présentent des ruptures de pente et éventuellement des paliers, pour les températures respectives $\theta_1 ~ [\si{\celsius}]$ et $\theta_2 ~ [\si{\celsius}]$ suivantes, la composition des systèmes étant donnée en fraction massique de plomb $w_{\ce{Pb}}$.

\begin{center}
\begin{tabular}{cccccccccccc}
$w_{\ce{Pb}}$ & 0.0 & 0.1 & 0.2 & 0.3 & 0.4 & 0.5 & 0.6 & 0.7 & 0.8 & 0.9 & 1.0 \\
$\theta_1 ~ [\si{\celsius}]$ & 321 & 316 & 308 & 302 & 293 & 284 & 274 & 263 & 251 & 272 & 288 \\
$\theta_2 ~ [\si{\celsius}]$ & $-$ & 250 & 250 & 250 & 250 & 250 & 250 & 250 & 250 & 250 & $-$ \\
\end{tabular}
\end{center}


\Q Tracer l'allure du diagramme binaire isobare cadmium-plomb en admettant que ces deux métaux ont une miscibilité nulle à l'état solide. Identifier les différents domaines de ce diagramme.

\Q Évaluer les coordonnées du point eutectique, noté \textbf{E}.

\Q Une masse $m_1 = 1.00 ~ \si{\kilo\gram}$ d'un mélange liquide cadmium-plomb, équimolaire en chacun des deux métaux, est refroidie très progressivement.

\begin{enumerate}
\item Calculer, en fraction massique, la composition $w_{\ce{Pb}}$ du mélange initial.

\item Donner l'ordre de grandeur de la température de début de solidification et la nature du solide qui apparaît le premier.

\item Déterminer la masse maximale de ce métal qu'il est possible de cristalliser sans que l'autre métal ne cristallise.
\end{enumerate}

\Q De manière isotherme, du cadmium est ajouté à une masse $m_2 = 100 ~ \si{\gram}$ d'un mélange représenté par le point \textbf{M} de coordonnées $w_{\ce{Pb}} = 0.50$ et $\theta = 302 ~ \si{\celsius}$.

\begin{enumerate}
\item Préciser la composition $w^l_{\ce{Pb}}$, du liquide lorsque le premier cristal apparaît.

\item Déterminer la masse $\Delta m$ de cadmium pur qu'il a fallu  ajouter pour observer le début de cristallisation.
\end{enumerate}

\paragraph*{Données} 
\begin{align*}
M ( \ce{Cd} ) &= 112 ~ \si{\gram\per\mole} ; \\
M ( \ce{Pb} ) &= 207 ~ \si{\gram\per\mole} .
\end{align*}
